% Point 2: correlation, label encodings, and proper scores for LLM judges.
% Note answering the "evaluation" point of an advisor call (27 Nov 2026).
% Compile: latexmk -pdf point2_encodings_vs_first_moments_27Nov2026.tex
\documentclass[11pt]{article}
\usepackage[margin=1in]{geometry}
\usepackage{amsmath,amssymb,amsthm}
\usepackage{booktabs,array}
\usepackage[dvipsnames]{xcolor}
\usepackage[colorlinks=true,linkcolor=NavyBlue,urlcolor=NavyBlue,citecolor=NavyBlue]{hyperref}
\setlength{\parskip}{0.45em}
\setlength{\parindent}{0pt}
\sloppy\setlength{\emergencystretch}{3em}

\newtheorem{proposition}{Proposition}
\newtheorem{corollary}[proposition]{Corollary}
\theoremstyle{remark}
\newtheorem{remark}{Remark}
\theoremstyle{definition}
\newtheorem{example}{Example}

\newcommand{\E}{\mathbb{E}}
\newcommand{\Var}{\operatorname{Var}}
\newcommand{\Cov}{\operatorname{Cov}}
\newcommand{\Corr}{\operatorname{Corr}}
\newcommand{\pstar}{p^{*}}
\newcommand{\onehot}{e}
\newcommand{\RPS}{\mathrm{RPS}}
\newcommand{\Brier}{\mathrm{Brier}}
\newcommand{\CE}{\mathrm{CE}}
\newcommand{\1}{\mathbf{1}}
\newcommand{\file}[1]{\nolinkurl{#1}}
\newcolumntype{P}[1]{>{\raggedright\arraybackslash}p{#1}}

\newcommand{\tldr}[1]{\par\medskip\noindent\fbox{\parbox{\dimexpr\linewidth-2\fboxsep-2\fboxrule\relax}{\textbf{TL;DR.} #1}}\par\medskip}

\title{Point 2: Correlation, Encodings, and Proper Scores for LLM Judges}
\author{Brando Miranda}
\date{27 November 2026}

\begin{document}
\maketitle

\tldr{Pearson correlation between a judge's score and the human rating (i) ignores any affine distortion of the score, (ii) equals $r^2 = 1-\min_{a,b}\E(Y-a-bS)^2/\Var(Y)$, i.e.\ squared error \emph{after} a best affine recalibration fitted on the evaluation labels themselves, so it keeps (the linear part of) discrimination and gets calibration for free, and (iii) sees only the conditional mean of a scalar score, which does not determine the rating distribution when there are three or more categories. The fix is to choose the \emph{label encoding} first. Squared error on the one-hot encoding is the multiclass Brier score; squared error on the cumulative encoding $\1[Y\le k]$ is the ranked probability score (RPS). Both are strictly proper, RPS also respects the ordering, and each recovers $\pstar(\cdot\mid c)$ because $\E[\text{encoding}\mid c]$ does. A scalar encoding has rank one and recovers only the mean. Worked example below: a point mass at $2$ and a half--half split on $\{0,4\}$ have the same mean and tie on scalar squared error, but differ by $3/2$ in Brier and $1$ in RPS (normalized: $1/5$).}

\section{What the point says, and the one refinement}
\label{sec:point}

The point (my own write-up of the advisor discussion) has five parts, all of which hold:
\begin{enumerate}
\item Pearson correlation is invariant to affine maps, so a judge that scores every item one point too high still has $r=1$ (Section~\ref{sec:affine}).
\item $r^{2}=1-\mathrm{MSE}(\text{best affine fit})/\Var(Y)$, so correlation is squared error after an optimal linear recalibration; it captures discrimination and ignores calibration (Section~\ref{sec:identity}).
\item With $Y$ one-hot, $\E[Y\mid x]$ is the distribution and squared error on it is the multiclass Brier score, which is proper (Section~\ref{sec:onehot}).
\item Judge correlations are computed on a scalar 0--5 score, whose mean cannot recover the spread (Section~\ref{sec:scalar}).
\item RPS is the same squared error on the cumulative encoding $\1[Y\le k]$, so it keeps the full distribution and adds the ordering (Section~\ref{sec:cumulative}). The paper should be organised around encodings, not first moments (Section~\ref{sec:framing}).
\end{enumerate}

\textbf{Refinement of ``still improper''.} Correlation is a squared-error proxy only after an affine refit, and the sentence ``it can be a squared-error proxy and still be improper'' needs a precise scope. Correlation is not a per-item scoring rule at all, so ``proper'' has no direct meaning for it; the natural population-level readings give two different verdicts (Propositions~\ref{prop:expected} and~\ref{prop:sampled}). If the judge's scalar score is a deterministic function of the item (for example the mean of its forecast), the truthful forecast maximizes $r$ but \emph{not uniquely}: correlation is non-identifying, at best weakly proper. If the score is a sample drawn from the judge's forecast, correlation is strictly \emph{anti}-proper in the sense that a collapsed judge beats the truthful one. I use ``improper'' only for the second case and ``non-identifying'' for the first.

\section{Setup}
\label{sec:setup}

An item (prompt, response, rubric context) is $c$. The human rating is $Y\in\{0,1,\dots,K-1\}$ (for a 0--5 rubric, $K=6$). Conditional on $c$, the rater pool and the rubric, the rating has distribution
\[
\pstar(k\mid c)=\Pr(Y=k\mid C=c),\qquad k=0,\dots,K-1 .
\]
A \emph{judge forecast} is a probability vector $q(\cdot\mid c)\in\Delta^{K-1}$. Two scalar readouts matter:
\begin{itemize}
\item the \emph{expected-score readout} $s(c)=\mu_q(c):=\sum_k k\,q(k\mid c)$, a deterministic function of $c$;
\item the \emph{sampled readout} $S\sim q(\cdot\mid c)$, drawn independently of $Y$ given $C=c$ (how a judge that outputs one sampled 0--5 grade behaves).
\end{itemize}
Write $\mu(c)=\E[Y\mid C=c]=\sum_k k\,\pstar(k\mid c)$. All moments are over the joint law of $(C,Y)$ (and of $S$ where relevant). Pearson correlation is $r(S,Y)=\Cov(S,Y)/\sqrt{\Var(S)\Var(Y)}$, defined when both variances are positive.

For a forecast $q$ and a realized rating $y$, with $\onehot_y\in\{0,1\}^K$ the one-hot vector and $F_q(t)=\sum_{k\le t}q(k)$:
\begin{align*}
\Brier(q,y)&=\|q-\onehot_y\|_2^{2}=\sum_{k}\bigl(q(k)-\1[y=k]\bigr)^{2},\\
\RPS(q,y)&=\sum_{t=0}^{K-2}\bigl(F_q(t)-\1[y\le t]\bigr)^{2},\qquad \text{normalized: } \RPS/(K-1)\in[0,1],\\
\CE(q,y)&=-\ln q(y)\quad(\text{nats; }+\infty\text{ if }q(y)=0).
\end{align*}
These are the repository's conventions (\file{experiments/36\_.../metrics.py}: \texttt{rps\_one}, \texttt{normalized\_rps\_one}). A score $S(q,y)$ is \emph{proper} if $\E_{Y\sim p}S(p,Y)\le\E_{Y\sim p}S(q,Y)$ for all $p,q$ (loss orientation), and \emph{strictly proper} if equality holds only for $q=p$ (Gneiting and Raftery, 2007).

\section{Affine invariance}
\label{sec:affine}

\begin{proposition}[Affine invariance]\label{prop:affine}
For $a\neq0$ and any $b$, $r(aS+b,\,Y)=\operatorname{sgn}(a)\,r(S,Y)$.
\end{proposition}
\begin{proof}
$\Cov(aS+b,Y)=a\Cov(S,Y)$ and $\Var(aS+b)=a^{2}\Var(S)$, so
\[
r(aS+b,Y)=\frac{a\,\Cov(S,Y)}{|a|\sqrt{\Var S\,\Var Y}}=\operatorname{sgn}(a)\,r(S,Y).
\]
\end{proof}

\begin{example}[Shift and stretch]\label{ex:shiftstretch}
Let $S=Y+1$: every score is one point too high, yet $r(S,Y)=1$ while the mean squared error is $\E(S-Y)^2=1$. Let $S=2Y$: again $r=1$, and the error is $\E(S-Y)^2=\E[Y^2]$, which is large whenever the ratings are not concentrated near $0$ (on the eight-item list $0,1,2,3,4,5,2,3$ it is $17/2$). A correlation of one is compatible with a wrong location and a wrong spread.
\end{example}

\begin{remark}[Rank correlations]
Spearman's $\rho$ and Kendall's $\tau$ depend on the scores only through the pairwise order of items, so they are invariant under any strictly increasing transformation of $S$ (and flip sign under a strictly decreasing one). This is strictly stronger than affine invariance: $S=\exp(Y)$ also gets $\rho=\tau=1$. The affine-invariance statement for Pearson is Lean-checked in the repository for $a>0$ and rational data (\texttt{pearson\_affine\_invariant}, Section~\ref{sec:scope}); the rank-invariance statement is a one-line argument that I have not seen formalized.
\end{remark}

\section{The identity \texorpdfstring{$r^{2}=1-\min_{a,b}\mathrm{MSE}/\Var(Y)$}{r-squared = 1 - min MSE / Var(Y)} and what it keeps}
\label{sec:identity}

\begin{proposition}[Correlation is squared error after the best affine recalibration]\label{prop:identity}
If $\Var(S)>0$ and $\Var(Y)>0$, then
\[
\min_{a,b\in\mathbb{R}}\E\bigl[(Y-a-bS)^{2}\bigr]=\Var(Y)\bigl(1-r(S,Y)^{2}\bigr),
\qquad\text{i.e.}\qquad
r^{2}=1-\frac{\min_{a,b}\E(Y-a-bS)^{2}}{\Var(Y)},
\]
attained at $b^{*}=\Cov(S,Y)/\Var(S)$ and $a^{*}=\E Y-b^{*}\E S$.
\end{proposition}
\begin{proof}
For fixed $b$, $\E(Y-a-bS)^{2}=\Var(Y-bS)+(\E Y-b\E S-a)^{2}$ is minimized over $a$ by $a=\E Y-b\E S$ (the normal equation for the intercept), leaving $\Var(Y-bS)=\Var(Y)-2b\Cov(S,Y)+b^{2}\Var(S)$. This quadratic in $b$ is minimized at $b^{*}=\Cov(S,Y)/\Var(S)$ (the normal equation for the slope), with minimum value $\Var(Y)-\Cov(S,Y)^{2}/\Var(S)=\Var(Y)\bigl(1-r^{2}\bigr)$.
\end{proof}

The same algebra holds for the empirical (in-sample) moments of a finite evaluation set; the Python check \file{point2\_worked\_examples\_27Nov2026.py} verifies the identity in exact rational arithmetic on 200 random datasets.

\textbf{Reading.} The quantity $\Var(Y)\,r^{2}$ is what a squared-error score would give a judge \emph{after} the scores have been passed through the best affine map, a map fitted on the very labels being evaluated. Equivalently $r^{2}$ is the mean-squared-error skill score of the recalibrated score against the constant forecast $\E Y$. A judge whose scale or offset is wrong is not penalized: that is the content of Proposition~\ref{prop:affine}, and Proposition~\ref{prop:identity} says why.

\subsection{Link to the Murphy decomposition}

Let $g(s)=\E[Y\mid S=s]$ denote the calibration curve of the score.

\begin{proposition}[Uncertainty, resolution, reliability for squared error]\label{prop:murphy}
For any real score $S$ with finite second moment,
\[
\E(Y-S)^{2}\;=\;\underbrace{\Var(Y)}_{\text{uncertainty}}\;-\;\underbrace{\Var\bigl(g(S)\bigr)}_{\text{resolution}}\;+\;\underbrace{\E\bigl(g(S)-S\bigr)^{2}}_{\text{reliability (calibration error)}}.
\]
Let $S'=a^{*}+b^{*}S$ be the best affine recalibration (with $b^*\ne 0$). Then $g(S')=g(S)$ as a function of the information in the score, and
\[
r(S,Y)^{2}\,\Var(Y)\;=\;\Var\bigl(g(S)\bigr)\;-\;\E\bigl(g(S)-S'\bigr)^{2}.
\]
In words, $r^{2}\Var(Y)$ is the resolution minus the calibration error that remains \emph{after} the affine fix; it equals the resolution exactly when the calibration curve $g$ is affine.
\end{proposition}
\begin{proof}
$Y-g(S)$ is orthogonal to every square-integrable function of $S$, so $\E(Y-S)^{2}=\E(Y-g(S))^{2}+\E(g(S)-S)^{2}$, and $\E(Y-g(S))^{2}=\Var(Y)-\Var(g(S))$ by the law of total variance. For the second display, apply the same decomposition to $S'$, which carries the same information as $S$ (so $g(S')=g(S)$): $\E(Y-S')^{2}=\Var(Y)-\Var(g(S))+\E(g(S)-S')^{2}$. By Proposition~\ref{prop:identity} the left side equals $\Var(Y)(1-r^{2})$.
\end{proof}

For a binary outcome and a probability forecast this is the Brier-score decomposition of Murphy (1973), written there for forecasts grouped into bins; Br\"ocker (2009) extends it to any strictly proper score and finite-valued targets. Correlation therefore keeps resolution (discrimination) and discards the \emph{affine part} of reliability, the part an offset and a scale can repair. What remains of the reliability term after the fix is a nonlinear calibration error that correlation also does not charge for, beyond its effect on $r^{2}$.

\section{Why a squared-error proxy can still fail to be proper}
\label{sec:improper}

Three separate reasons, in decreasing order of generality.

\emph{(1) Correlation is not a per-item scoring rule.} Propriety is a property of $S(q,y)$: a function of one forecast and one outcome, averaged over items. $r$ is a ratio of population moments, and the affine refit $(a^{*},b^{*})$ is computed from all the labels. It cannot be written as $\E\,S(q(C),Y)$. Standard propriety results (Gneiting and Raftery, 2007) do not apply to it directly, so one has to ask the population-level question of which forecasts maximize it.

\emph{(2) The recalibration uses the evaluation labels.} The refit supplies calibration for free, so the score no longer rewards the judge for being calibrated. If $(a,b)$ were instead fitted on separate data, the refit would become an explicit, visible calibration step (this is the ``matched'' arm of Experiment 41, Section~\ref{sec:limits}).

\emph{(3) What it recovers is one moment.} Squared error on a scalar is strictly consistent for the \emph{mean} of the outcome distribution (Gneiting, 2011), nothing more (Section~\ref{sec:scalar}).

The next two propositions turn (1) into theorems for the two readouts.

\begin{proposition}[Expected-score readout: the truth maximizes $r$, but not uniquely]\label{prop:expected}
Let $s=s(C)$ be a function of the item alone with $\Var(s)>0$, and suppose $\Var(\mu(C))>0$. Then
\[
r(s,Y)\;=\;r\bigl(s,\mu(C)\bigr)\sqrt{\frac{\Var(\mu(C))}{\Var(Y)}}\;\le\;\sqrt{\frac{\Var(\mu(C))}{\Var(Y)}}=r\bigl(\mu(C),Y\bigr),
\]
with equality if and only if $s=a\mu(C)+b$ almost surely for some $a>0$ and $b$.
\end{proposition}
\begin{proof}
Since $s$ is a function of $C$, $\E[sY]=\E\bigl[s\,\E[Y\mid C]\bigr]=\E[s\mu]$, so $\Cov(s,Y)=\Cov(s,\mu)$. Also $\E Y=\E\mu$. Hence $r(s,Y)=\Cov(s,\mu)/\sqrt{\Var s\,\Var Y}=r(s,\mu)\sqrt{\Var\mu/\Var Y}$, and $r(s,\mu)\le1$ with equality iff $s$ is a positive affine function of $\mu$ (Cauchy--Schwarz).
\end{proof}

\begin{corollary}[Same conditional mean $\Rightarrow$ tie]\label{cor:tie}
Two forecasts $q,q'$ with $\mu_q(c)=\mu_{q'}(c)$ for all $c$ have the same expected-score readout, hence the same Pearson, Spearman and Kendall correlation with $Y$ and the same mean squared error of that readout, whatever their spread. Any forecast whose mean is a positive affine image of $\mu(c)$ attains the maximal $r$ of Proposition~\ref{prop:expected}, including the truthful $\pstar$ and, when $\mu(c)$ lies on the rating grid, the point mass at $\mu(c)$.
\end{corollary}
(This is the ``same-mean non-identification'' statement. The tie is a one-line consequence of all four statistics being functions of the readout alone; the repository's Lean files check concrete same-mean pairs, not the general statement, see Section~\ref{sec:scope}.)

\begin{proposition}[Sampled readout: correlation is anti-proper]\label{prop:sampled}
Let the judge's score be $S\sim q(\cdot\mid C)$ independent of $Y$ given $C$. Put $\rho=\Var(\mu(C))/\Var(Y)$ and suppose $\Var(\mu(C))>0$ and $\E\Var(Y\mid C)>0$.
\begin{enumerate}
\item The truthful judge $q=\pstar$ has $r(S,Y)=\rho$, and $0<\rho<1$.
\item The collapsed judge that outputs the deterministic score $S_\delta=a\mu(C)+b$ with $a>0$ has $r(S_\delta,Y)=\sqrt{\rho}>\rho$.
\end{enumerate}
So correlation ranks the collapsed judge strictly above the truthful one. When $\mu(c)$ is on the rating grid and $a=1,b=0$, $S_\delta$ is the point-mass forecast $\delta_{\mu(c)}$; by strict propriety of Brier, RPS and CE it is strictly worse than $\pstar$ on each of them whenever $\pstar(\cdot\mid c)$ is not a point mass on a set of items of positive probability, which $\E\Var(Y\mid C)>0$ guarantees.
\end{proposition}
\begin{proof}
(1) $\E[S\mid C]=\mu(C)$, so $\E[SY]=\E[\mu^{2}]$, $\E S=\E Y$ and $\Cov(S,Y)=\E\mu^{2}-(\E Y)^{2}=\Var\mu$. The marginal law of $S$ equals that of $Y$, so $\Var S=\Var Y$. Thus $r=\Var\mu/\Var Y=\rho$. By the law of total variance $\Var Y=\Var\mu+\E\Var(Y\mid C)$, so $0<\rho<1$.
(2) $\Cov(S_\delta,Y)=a\Var\mu$ and $\Var S_\delta=a^{2}\Var\mu$, so $r=a\Var\mu/\bigl(a\sqrt{\Var\mu\,\Var Y}\bigr)=\sqrt{\rho}$, which exceeds $\rho$ because $0<\rho<1$.
\end{proof}

\begin{remark}
Proposition~\ref{prop:sampled} is the correlation analogue of the familiar fact that hedging to the mean beats honest sampling under a loss that sees only first moments; it is the reason the paper should say ``correlation is anti-proper for a sampled readout'' only with that qualifier, and ``non-identifying'' for the expected-score readout (Proposition~\ref{prop:expected}).
\end{remark}

\subsection{Scope of the formal checks in the repository}
\label{sec:scope}

The repository's standard (AGENTS.md) is that a claim is called a theorem only if it is checked in Lean~4 and in Python, and that the squared form of anti-properness is the Lean-checked one. Here is exactly what I found, and what I did not.

\begin{center}
\small
\begin{tabular}{P{4.0cm}P{5.2cm}P{5.2cm}}
\toprule
Claim in this note & Repository evidence & Status \\
\midrule
Brier is strictly proper (any $K$) & \file{experiments/35\_brier\_vs\_corr\_math\_l2\_pearson\_lean\_start\_26Nov2026/BrierCorr.lean}: \texttt{brier\_strictly\_proper}, \texttt{brier\_decomposition} & Lean, rationals, standard axioms only (\texttt{lean\_check.txt}) \\
Pearson affine invariance (Prop.~\ref{prop:affine}) & same file: \texttt{pearson\_affine\_invariant} ($a>0$); also \texttt{mse\_decomposition}, \texttt{standardized\_mse\_corr} & Lean; hypotheses supply rational standard deviations; $a<0$ and Spearman/Kendall are paper arguments \\
Same-mean pairs differ in Brier/RPS & \file{experiments/35\_.../extended\_checks\_11-26-2026/squared\_difference\_vs\_distribution\_11-26-2026/MeanMatching.lean}; the Lean file under \file{experiments/36\_.../}: \texttt{ordinal\_same\_mean\_different\_law} & Lean, but the second is arithmetic on typed numbers, not linked to a Lean definition of RPS (the repository's own QA says so). General Cor.~\ref{cor:tie} is a paper argument \\
Worked-example table (Section~\ref{sec:scalar}) & new: \file{Point2WorkedExamples.lean} defines Brier and RPS on ratings $0..5$ and evaluates all entries & Lean (\texttt{decide +kernel}), no \texttt{sorry}, standard axioms; a deliberately wrong entry fails to compile \\
Anti-properness, squared form $\rho^{2}<\rho$ & the Lean file under \file{experiments/36\_.../}: \texttt{anti\_proper\_from\_parts} & Lean; proves only $0<\rho<1\Rightarrow\rho^{2}<\rho$ \\
Anti-properness, square-root form (Prop.~\ref{prop:sampled}) & AGENTS.md: Monte Carlo only. A finite-joint-law Lean theorem \texttt{Model.anti\_proper} with the strict inequality exists in \file{PlanClaims.lean} in the plan-claims audit folder under experiment~35's \texttt{extended\_checks\_11-26-2026/} & I did \emph{not} recompile it (needs Mathlib). Treat the square-root form as paper-proved plus exact-fraction Python (Part~2 of the script) until the AGENTS.md wording and that file are reconciled \\
Props.~\ref{prop:identity}, \ref{prop:murphy}, \ref{prop:expected}, \ref{prop:encoding}, \ref{prop:rps} & this note & Paper proofs; identity and RPS/Brier expectations also checked in exact fractions (200 random cases each); \emph{no new Lean} \\
\bottomrule
\end{tabular}
\end{center}

\section{One-hot encoding and the Brier score}
\label{sec:onehot}

\begin{proposition}[One-hot squared error is the multiclass Brier score, strictly proper]\label{prop:onehot}
$\E[\onehot_Y\mid C=c]=\pstar(\cdot\mid c)$, and for every forecast $q$,
\[
\E_{Y\sim\pstar}\,\Brier(q,Y)=\|q-\pstar\|_2^{2}+\bigl(1-\|\pstar\|_2^{2}\bigr).
\]
Hence Brier is strictly proper: the expected score is minimized uniquely at $q=\pstar$, with minimum $1-\|\pstar\|_2^{2}$.
\end{proposition}
\begin{proof}
The $k$th coordinate of $\onehot_Y$ is $\1[Y=k]$, with mean $\pstar(k)$. Then $\E\|q-\onehot_Y\|^{2}=\|q\|^{2}-2\langle q,\E\onehot_Y\rangle+\E\|\onehot_Y\|^{2}=\|q\|^{2}-2\langle q,\pstar\rangle+1=\|q-\pstar\|^{2}+1-\|\pstar\|^{2}$, using $\|\onehot_Y\|^{2}=1$.
\end{proof}

The logarithmic score is the other standard strictly proper score for a categorical label: $\E_{Y\sim\pstar}\CE(q,Y)=H(\pstar)+\mathrm{KL}(\pstar\,\|\,q)$ by Gibbs' inequality (Gneiting and Raftery, 2007). It is not a squared error on an encoding, it is blind to ordering, and it is infinite whenever $q$ puts zero mass on a realized rating.

\begin{remark}[$K=2$ is the one case where the scalar suffices]
For $K=2$ the mean is $\pstar(1\mid c)$, which \emph{is} the distribution, and $\Brier=2(q_1-y)^{2}$, twice the scalar squared error (the repository's \texttt{binary\_brier\_eq\_two\_l2}). The argument below therefore bites exactly when $K\ge3$, which includes every 0--5 rubric.
\end{remark}

\section{The scalar encoding sees one moment}
\label{sec:scalar}

Squared error on the scalar $Y$ is the encoding $\varphi(y)=y$, i.e.\ the row $A=(0,1,\dots,K-1)$ applied to $\onehot_y$. Its conditional expectation is $\mu(c)=\sum_k k\,\pstar(k\mid c)$, a single linear functional of $\pstar$. The general statement covers all three encodings of this note.

\begin{proposition}[Encoding rank]\label{prop:encoding}
Let $A\in\mathbb{R}^{m\times K}$ and define the encoded squared-error score $S_A(q,y)=\|A(q-\onehot_y)\|_2^{2}$. Then, for every $p$ and $q$,
\[
\E_{Y\sim p}S_A(q,Y)=\|A(q-p)\|_2^{2}+\operatorname{tr}\bigl(A\,\Sigma_p A^{\top}\bigr),\qquad \Sigma_p=\operatorname{diag}(p)-pp^{\top},
\]
so $S_A$ is proper. It is strictly proper if and only if $A$ is injective on $D=\{v\in\mathbb{R}^{K}:\sum_k v_k=0\}$. If some nonzero $v\in D$ has $Av=0$, then $q$ and $q+tv$ (for small $t$, staying in the simplex) have the same expected score under every truth.
\end{proposition}
\begin{proof}
$q-\onehot_Y=(q-p)-(\onehot_Y-p)$ and $\E[\onehot_Y-p]=0$, so $\E\|A(q-\onehot_Y)\|^{2}=\|A(q-p)\|^{2}+\E\|A(\onehot_Y-p)\|^{2}$, and the last term equals $\operatorname{tr}(A\Sigma_pA^{\top})$ and does not depend on $q$. The first term vanishes iff $A(q-p)=0$; since $q-p\in D$, this forces $q=p$ exactly when $\ker A\cap D=\{0\}$.
\end{proof}

\begin{itemize}
\item \textbf{One-hot} $A=I_K$: injective, strictly proper (Proposition~\ref{prop:onehot}).
\item \textbf{Cumulative} $A=L$ with rows $t=0,\dots,K-2$ and $(Lv)_t=\sum_{k\le t}v_k$: if $Lv=0$ then $v_0=\dots=v_{K-2}=0$, and $v\in D$ gives $v_{K-1}=0$. So $L$ is injective on $D$: strictly proper (Section~\ref{sec:cumulative}).
\item \textbf{Scalar} $A=(0,1,\dots,K-1)$: rank one, so $\ker A\cap D$ has dimension $K-2$. For $K\ge3$ it is nonzero, and every pair of laws with equal mean ties exactly. For $K=2$ it is $\{0\}$ and the scalar is faithful (previous remark). The scalar squared error is strictly consistent for the mean functional (Gneiting, 2011), which is a different property from being proper for the distribution.
\end{itemize}

Pearson correlation then discards a further two parameters (the affine map in Propositions~\ref{prop:affine} and~\ref{prop:identity}), and Spearman and Kendall discard all monotone maps.

\subsection{Worked example: same mean, different spread}

Take $K=6$ (ratings $0,\dots,5$) and two laws:
\[
A=\delta_{2}=(0,0,1,0,0,0),\qquad B=\tfrac12\delta_{0}+\tfrac12\delta_{4}=(\tfrac12,0,0,0,\tfrac12,0).
\]
Both have mean $2$; the variances are $0$ and $4$. Their cumulative distribution functions at $t=0,\dots,4$ are $F_A=(0,0,1,1,1)$ and $F_B=(\tfrac12,\tfrac12,\tfrac12,\tfrac12,1)$. A judge issuing forecast $A$ or $B$ has the same scalar readout $2$, so the scalar squared error of the readout is the same for both: $0$ if the truth is $A$ and $4$ if the truth is $B$. Using Proposition~\ref{prop:onehot} and its RPS analogue (Proposition~\ref{prop:rps}), the expected scores of forecast $q$ under truth $p$ are:

\begin{center}
\begin{tabular}{llcccc}
\toprule
 & & \multicolumn{2}{c}{truth $A$} & \multicolumn{2}{c}{truth $B$} \\
\cmidrule(lr){3-4}\cmidrule(lr){5-6}
Score & forecast $\to$ & $A$ & $B$ & $A$ & $B$ \\
\midrule
Scalar squared error of the mean readout & & $0$ & $0$ & $4$ & $4$ \\
Brier $\|q-\onehot_Y\|^{2}$ & & $0$ & $3/2$ & $2$ & $1/2$ \\
RPS (raw, thresholds $0..4$) & & $0$ & $1$ & $2$ & $1$ \\
RPS normalized by $K-1=5$ & & $0$ & $1/5$ & $2/5$ & $1/5$ \\
Cross-entropy (nats), exact & & $0$ & $+\infty$ & $+\infty$ & $\ln 2\approx0.6931$ \\
\bottomrule
\end{tabular}
\end{center}
Read across: under truth $A$, forecasts $A$ and $B$ score as in the first two columns; under truth $B$, as in the last two. Entries are expectations under the stated truth; lower is better. Brier and RPS values are exact fractions computed by \file{point2\_worked\_examples\_27Nov2026.py} and checked in Lean by \file{Point2WorkedExamples.lean}; the cross-entropy entries use exact zeros of $q$. (With 10\% uniform smoothing of the forecast, a numerical device and not part of the theory, the finite cross-entropies would be $0.0870$ and $0.7621$ on the diagonal and $4.0943$ off it, from the same script.)

\textbf{Reading.} The scalar score cannot tell the two forecasts apart whichever truth holds, and Corollary~\ref{cor:tie} says the same for Pearson, Spearman and Kendall. Brier and RPS separate them under both truths, with the truthful forecast strictly best, as Proposition~\ref{prop:onehot} and Proposition~\ref{prop:rps} require. The Brier gap between forecasts under truth $A$ is $\|B-A\|^{2}=3/2$ and the RPS gap is $1$, which the repository's \texttt{ordinal\_same\_mean\_different\_law} also records.

\subsection{A population version: correlation is maximal for a judge that is worse than the base rate}

Take three equally likely items with truths $\pstar(\cdot\mid c_1)=\tfrac12\delta_0+\tfrac12\delta_2$, $\pstar(\cdot\mid c_2)=\tfrac12\delta_0+\tfrac12\delta_4$ and $\pstar(\cdot\mid c_3)=\tfrac12\delta_2+\tfrac12\delta_4$. The conditional means are $1,2,3$, the conditional variances $1,4,1$, and $\Var(Y)=8/3=\Var\mu+\E\Var(Y\mid C)=2/3+2$. Compare the truthful judge $T$ ($q=\pstar$) with the collapsed judge $D$ ($q=\delta_{\mu(c)}$):

\begin{center}
\begin{tabular}{lcc}
\toprule
Population-average score (lower is better, except the last row) & $T$ & $D$ \\
\midrule
Pearson $r$, expected-score readout (identical $s=\mu$) & $1/2$ & $1/2$ \\
Pearson $r$, \emph{sampled} readout (Prop.~\ref{prop:sampled}) & $1/4$ & $1/2$ \\
Brier & $1/2$ & $2$ \\
RPS (raw) & $2/3$ & $4/3$ \\
RPS normalized ($/5$) & $2/15$ & $4/15$ \\
Cross-entropy (nats) & $\ln2\approx0.6931$ & $+\infty$ \\
Brier skill score $1-S(\text{judge})/S(\text{base rate})$ & $1/4$ & $-2$ \\
RPS skill score $1-S(\text{judge})/S(\text{base rate})$ & $1/4$ & $-1/2$ \\
\bottomrule
\end{tabular}
\end{center}
The base-rate forecast is the pooled marginal $(\tfrac13,0,\tfrac13,0,\tfrac13,0)$, with Brier $2/3$ and raw RPS $8/9$. On the expected-score readout the two judges tie on every correlation and on the mean squared error of the readout, and $r^{2}=1/4$ is the best any function of $c$ can achieve (Proposition~\ref{prop:expected}); yet $D$ is \emph{worse than the base rate} on Brier and RPS (negative skill) and has infinite cross-entropy. On the sampled readout, correlation prefers $D$ outright ($1/2$ against $1/4$), as Proposition~\ref{prop:sampled} predicts. These are constructed laws that verify the propositions; they are not evidence about any real judge.

\section{The cumulative encoding and the ranked probability score}
\label{sec:cumulative}

Let $Z_t=\1[Y\le t]$ for $t=0,\dots,K-2$ (the $t=K-1$ indicator is identically one). Then $\E[Z_t\mid C=c]=F^{*}(t\mid c)=\sum_{k\le t}\pstar(k\mid c)$, and $F^{*}(0),\dots,F^{*}(K-2)$ determine $\pstar$ (take successive differences, and $\pstar(K-1)=1-F^{*}(K-2)$). Squared error on this encoding is the ranked probability score (Epstein, 1969; Murphy, 1971):
\[
\RPS(q,y)=\sum_{t=0}^{K-2}\bigl(F_q(t)-\1[y\le t]\bigr)^{2}=\|L(q-\onehot_y)\|_2^{2},
\]
the sum over thresholds of the binary Brier score for the event $\{Y\le t\}$. The repository normalizes by $K-1$, so each threshold term lies in $[0,1]$ and $\RPS/(K-1)\in[0,1]$.

\begin{proposition}[RPS is strictly proper]\label{prop:rps}
$\displaystyle \E_{Y\sim p}\RPS(q,Y)=\sum_{t=0}^{K-2}\bigl(F_q(t)-F_p(t)\bigr)^{2}+\sum_{t=0}^{K-2}F_p(t)\bigl(1-F_p(t)\bigr)$, so the unique minimizer over $q$ is $q=p$.
\end{proposition}
\begin{proof}
For each $t$, $Z_t$ is Bernoulli with mean $F_p(t)$, and the binary Brier identity (Proposition~\ref{prop:onehot} with $K=2$) gives $\E(F_q(t)-Z_t)^{2}=(F_q(t)-F_p(t))^{2}+F_p(t)(1-F_p(t))$. Sum over $t$. The first sum is zero iff $F_q(t)=F_p(t)$ for all $t\le K-2$, iff $q=p$. (Equivalently, Proposition~\ref{prop:encoding} with $A=L$.)
\end{proof}

\begin{example}[Distance sensitivity: Brier is blind, RPS is not]\label{ex:distance}
Let the realized rating be $Y=0$, and compare the point-mass forecasts $\delta_{1}$ and $\delta_{5}$. Both put zero mass on the realized rating and the same mass elsewhere, so $\Brier=2$ for each (and the cross-entropy is infinite for each). For $\delta_j$ and outcome $y$, $F_q(t)=\1[j\le t]$ and $\1[y\le t]$ differ exactly for $t$ between $\min(j,y)$ and $\max(j,y)-1$, so $\RPS(\delta_j,y)=|j-y|$. Thus $\RPS(\delta_1,0)=1$ and $\RPS(\delta_5,0)=5$ (normalized $1/5$ and $1$). Brier charges every wrong category the same; RPS charges the distance. The same values are computed by the script and checked in \file{Point2WorkedExamples.lean} (\texttt{distance\_sensitivity}).
\end{example}

\begin{remark}[CRPS as the continuous limit]
For a real-valued outcome and a forecast cumulative distribution function $F$, $\mathrm{CRPS}(F,y)=\int_{-\infty}^{\infty}\bigl(F(t)-\1[y\le t]\bigr)^{2}dt$, the integral of the binary Brier scores at all thresholds (Matheson and Winkler, 1976; Gneiting and Raftery, 2007, Section~4.2); it is strictly proper on distributions with finite first moment. For a forecast supported on the integers and an integer outcome, $F$ and $\1[y\le t]$ are constant on each $[k,k+1)$, so the CRPS equals the unnormalized RPS: the RPS is the CRPS on the integer grid, and continuous ratings are handled by the same recipe with the cumulative encoding $\{\1[y\le t]\}_{t\in\mathbb{R}}$.
\end{remark}

\section{Framing recommendation: encoding, then matched strictly proper score}
\label{sec:framing}

\begin{center}
\small
\begin{tabular}{P{2.9cm}P{3.0cm}P{3.1cm}P{3.0cm}P{2.4cm}}
\toprule
Encoding $\varphi(Y)$ & $\E[\varphi\mid c]$ recovers & Squared-error score on $\varphi$ & Proper for $\pstar$? & Ordering-aware? \\
\midrule
Scalar $Y$ & the mean $\mu(c)$ only (rank-one) & $(s-y)^{2}$ & consistent for the mean; not proper for $\pstar$ when $K\ge3$ & uses numeric distance, spread invisible \\
Scalar, then affine refit (Pearson) & $\mu(c)$ up to an affine map & $1-r^{2}$ (in-sample refit) & no: non-identifying (Prop.~\ref{prop:expected}), anti-proper if sampled (Prop.~\ref{prop:sampled}) & only through the means \\
One-hot $\onehot_Y$ & $\pstar(\cdot\mid c)$ & Brier $\|q-\onehot_y\|^{2}$ & strictly & no: every wrong category costs the same \\
Cumulative $\1[Y\le t]$, $t\le K-2$ & $F^{*}(\cdot\mid c)$, hence $\pstar$ & RPS & strictly & yes \\
$\1[Y\le t]$, all real $t$ & $F^{*}(t\mid c)$, $t\in\mathbb{R}$ & CRPS & strictly (finite first moment) & yes \\
\bottomrule
\end{tabular}
\end{center}
Cross-entropy (the log score) on the one-hot encoding is a fifth strictly proper option, not a squared error; it is ordering-blind and infinite on zero-mass outcomes.

\textbf{Recommendation.} State the validation target as the forecast of $\pstar(y\mid c,\text{rater pool},\text{rubric})$ and organize the evaluation protocol around the label type:
\begin{enumerate}
\item \emph{Declare the encoding by label type}: categorical $\to$ one-hot (CE or Brier); ordinal $\to$ cumulative (RPS); continuous $\to$ CRPS.
\item \emph{Score with the matched strictly proper score}, as a per-item held-out mean, with item-clustered bootstrap intervals (the repository's standing protocol).
\item \emph{Report the base-rate skill score} $1-S(\text{judge})/S(\text{base rate})$, so that a number is read against what a judge that ignores the item would get (the last two rows of the population table show why this matters).
\item \emph{Stratify by human agreement}, since the skill ceiling depends on how peaked $\pstar$ is.
\item \emph{Keep correlation as a discrimination diagnostic} next to a calibration diagnostic, not as the validity criterion.
\end{enumerate}
This replaces a framing around first moments (mean, variance, correlation) by one around what $\E[\text{encoding}\mid c]$ identifies. It is also where the paper's novelty sits: neither ``use proper scores'' nor ``labels are distributions'' is new; what is specific is the practice gap in judge validation, the identification result for correlation, the encoding-matched protocol, and the evidence on the size of the gap.

\section{What this does and does not establish}
\label{sec:limits}

\begin{itemize}
\item \textbf{Correlation remains useful.} By Proposition~\ref{prop:murphy}, $r^{2}\Var(Y)$ is a (linear) resolution term. If a downstream use only needs the ordering of items (best-of-$n$ selection, ranking), affine and monotone invariance is a feature. The claim here is about validating a judge as a \emph{forecaster of the human rating distribution}.
\item \textbf{The constructed examples are not real judges.} Propositions~\ref{prop:expected} and~\ref{prop:sampled} and the worked examples show that the metrics \emph{can} disagree, exactly. A rank reversal between real LLM judges under a pre-registered paired test is not established by this note.
\item \textbf{Exp.~41: calibration closes most of the gap.} Experiment~41 (\file{experiments/41\_output\_space\_k\_sweep\_g\_d\_study\_sufficiency\_and\_label\_dollars\_start\_27Nov2026/results/k\_sweep\_summary\_table.md}) trains \emph{linear-head} judges (not LLM judges) with cross-entropy or with Pearson correlation on seven real multi-annotator rating sets, viewed at $K\in\{2,3,5\}$ (17 views; ChaosNLI and $\alpha$NLI have one native view each), at a fixed ridge penalty of $0.01$ and at most 35 optimizer iterations. Held-out normalized RPS, mean over five split seeds, ``Pearson-trained minus CE-trained'' (positive means the CE-trained head is better): the raw gap is up to $+0.1282$ (largest in the SummEval crowd $K=5$ view, which the table flags as carrying no signal: constant-forecast normalized RPS $0.1277$); after each head's score passes through its own validation-fitted multinomial-logit calibration map, the gap is at most $+0.0024$ (largest in POPQUORN politeness, $K=2$; single-seed values in that view range from $+0.0019$ to $+0.0031$). These are the table's five-seed means; I make no significance claim from them here. This is what Proposition~\ref{prop:murphy} predicts qualitatively: Pearson training buys discrimination, and the missing ingredient is calibration, which a post-hoc map supplies. It is a statement about these linear heads, not about frontier LLM judges.
\end{itemize}

\section{Suggested paragraph for the paper}
\label{sec:paragraph}

\begin{quote}
\small
\begin{verbatim}
\paragraph{Why correlation cannot validate a judge as a forecaster.}
Pearson correlation is invariant to affine maps of the judge's score, and
$r^{2}=1-\min_{a,b}\mathbb{E}(Y-a-bS)^{2}/\mathrm{Var}(Y)$ is the squared
error after the best linear recalibration fitted on the evaluation labels
themselves. It therefore rewards discrimination and does not charge for
miscalibration. A scalar score also cannot identify the rating distribution:
its conditional mean is one moment, so two forecasts with the same mean but
different spread tie on correlation and on squared error of the mean, while
differing in Brier, RPS and cross-entropy. Encoding the label first removes
the problem. With the one-hot encoding, $\mathbb{E}[e_Y\mid c]=p^{*}(\cdot\mid c)$
and squared error is the multiclass Brier score, strictly proper; with the
cumulative encoding $\mathbf{1}[Y\le k]$ it is the ranked probability score,
strictly proper and sensitive to the ordering of the categories. We therefore
validate a judge by the strictly proper score matched to the label encoding,
report its base-rate skill score stratified by human agreement, and keep
correlation as a diagnostic of discrimination.
\end{verbatim}
\end{quote}

\section{Verification ledger and references}
\label{sec:ledger}

\textbf{Checks run for this note} (files in \file{docs/notes/point2\_support\_27Nov2026/}): the exact-fraction script (\file{point2\_worked\_examples\_27Nov2026.py}, output \file{point2\_worked\_examples\_27Nov2026\_output.txt}, ends with \texttt{ALL ASSERTIONS PASSED}); the Lean file \file{Point2WorkedExamples.lean} with its compiler output \file{lean\_check.txt} (Lean 4.36.0, no Mathlib, no \texttt{sorry}, axioms \texttt{propext}, \texttt{Classical.choice}, \texttt{Quot.sound}). Reproduce: \texttt{python3 point2\_worked\_examples\_27Nov2026.py} and \texttt{lean Point2WorkedExamples.lean}.

\textbf{Not Lean-checked here:} Propositions~\ref{prop:identity}, \ref{prop:murphy}, \ref{prop:expected}, \ref{prop:sampled}, \ref{prop:encoding} and~\ref{prop:rps} are proved on paper above, with the numerical checks listed. Under the repository's three-pillar standard they should be called ``analytical'' until a zero-\texttt{sorry} Lean lemma exists.

\begin{thebibliography}{99}\small
\bibitem{brier1950} G.~W. Brier. Verification of forecasts expressed in terms of probability. \emph{Monthly Weather Review}, 78(1):1--3, 1950.
\bibitem{epstein1969} E.~S. Epstein. A scoring system for probability forecasts of ranked categories. \emph{Journal of Applied Meteorology}, 8:985--987, 1969.
\bibitem{murphy1971} A.~H. Murphy. A note on the ranked probability score. \emph{Journal of Applied Meteorology}, 10:155--156, 1971.
\bibitem{murphy1973} A.~H. Murphy. A new vector partition of the probability score. \emph{Journal of Applied Meteorology}, 12(4):595--600, 1973.
\bibitem{matheson1976} J.~E. Matheson and R.~L. Winkler. Scoring rules for continuous probability distributions. \emph{Management Science}, 22(10):1087--1096, 1976.
\bibitem{degroot1983} M.~H. DeGroot and S.~E. Fienberg. The comparison and evaluation of forecasters. \emph{The Statistician}, 32(1/2):12--22, 1983.
\bibitem{gr2007} T.~Gneiting and A.~E. Raftery. Strictly proper scoring rules, prediction, and estimation. \emph{Journal of the American Statistical Association}, 102(477):359--378, 2007.
\bibitem{brocker2009} J.~Br\"ocker. Reliability, sufficiency, and the decomposition of proper scores. \emph{Quarterly Journal of the Royal Meteorological Society}, 135(643):1512--1519, 2009.
\bibitem{gneiting2011} T.~Gneiting. Making and evaluating point forecasts. \emph{Journal of the American Statistical Association}, 106(494):746--762, 2011.
\end{thebibliography}

Citation use in the text: Brier (1950), the quadratic score; Epstein (1969) and Murphy (1971), the RPS; Murphy (1973) and Br\"ocker (2009), the uncertainty--resolution--reliability decomposition (DeGroot and Fienberg, 1983, for the calibration/refinement form of the same idea); Matheson and Winkler (1976) and Gneiting and Raftery (2007), the CRPS as an integral of threshold Brier scores and propriety; Gneiting (2011), squared error as a consistent scoring function for the mean. I confirmed each reference's venue, volume and pages against the publisher, DOI or author-hosted record, and the CRPS-as-integral-of-Brier statement against the text of Gneiting and Raftery (2007), Section~4.2, Eq.~(20). I did not re-derive the claims those papers make beyond what is proved above.

\tldr{Correlation is invariant to affine maps and equals squared error after a best affine refit on the evaluation labels, so it measures discrimination and gives calibration away; for a deterministic readout the truth ties with every positively affine image of its mean (non-identifying), and for a sampled readout a collapsed judge beats the truth (anti-proper). A scalar encoding has rank one and sees only the mean; one-hot (Brier) and cumulative (RPS) encodings are injective on the simplex, so squared error on them is strictly proper and RPS also charges distance. Worked example: same mean $2$, point mass at $2$ versus half on $0$ and half on $4$: scalar error ties, Brier differs by $3/2$ and RPS by $1$ ($1/5$ normalized). Organize the paper around label encoding $\to$ matched strictly proper score $\to$ base-rate skill score stratified by human agreement, keep correlation as a diagnostic, and note that Experiment~41's linear heads show most of the raw gap (up to $0.1282$) vanishes after post-hoc calibration (at most $0.0024$); real-judge rank reversal remains unestablished.}

\end{document}
